\documentclass[10pt]{article}
\usepackage[T1]{fontenc}
\usepackage[utf8]{inputenc}
\usepackage[french]{babel}
\usepackage{amsmath,amssymb,mathrsfs,mathpazo}
\usepackage{geometry}
\geometry{a4paper,margin=2cm}
\usepackage{graphicx,array,tabularx,tabularray,colortbl,multirow,multicol,booktabs,enumitem,pifont,ulem,fancybox,diagbox,cancel}
\usepackage{pstricks,pstricks-add,pst-plot,pst-tree,pst-node,pst-text,pst-func,pst-math,pst-bspline,pst-3dplot,pst-eucl,pst-solides3d}
\let\pscalculate\undefined
\usepackage{pst-bezier}
\usepackage{tikz,pgfplots,tkz-tab}
\usetikzlibrary{arrows,arrows.meta,patterns,calc,positioning,trees,decorations.markings,shapes.geometric,shapes.misc,shapes.arrows,intersections,decorations.pathmorphing,decorations.pathreplacing,backgrounds,shadows,babel}
\pgfplotsset{compat=1.18}
\usepgfplotslibrary{fillbetween}
\usepackage[np]{numprint}
\usepackage{listings,makecell,pas-tableur,twoopt}
\let\barre\undefined
\usepackage[pstricks]{tablvar}
\definecolor{darkgreen}{RGB}{0,100,0}
\providecommand{\red}{\color{red}}
\providecommand{\blue}{\color{blue}}
\providecommand{\green}{\color{green}}
\providecommand{\black}{\color{black}}
\providecommand{\white}{\color{white}}
\providecommand{\gray}{\color{gray}}
\providecommand{\cyan}{\color{cyan}}
\providecommand{\magenta}{\color{magenta}}
\providecommand{\yellow}{\color{yellow}}
\providecommand{\violet}{\color{violet}}
\providecommand{\orange}{\color{orange}}
\providecommand{\euro}{\texteuro}
\providecommand{\EUR}{\texteuro}
\providecommand{\ieme}{\textsuperscript{e}}
\providecommand{\degres}{\ensuremath{^\circ}}
\providecommand{\degree}{\ensuremath{^\circ}}
\providecommand{\vv}[1]{\overrightarrow{#1}}
\providecommand{\Coord}[1]{#1}
\providecommand{\decosix}{}
\providecommand{\starredbullet}{\ensuremath{\bullet}}
\pagestyle{empty}
\setlength{\parindent}{0pt}
\usepackage{adjustbox}\DeclareUnicodeCharacter{2236}{:}\DeclareUnicodeCharacter{2190}{\ensuremath{\leftarrow}}
\providecommand{\R}{\mathbb{R}}
\providecommand{\N}{\mathbb{N}}
\providecommand{\Z}{\mathbb{Z}}
\providecommand{\Q}{\mathbb{Q}}
\providecommand{\D}{\mathbb{D}}
\providecommand{\C}{\mathbb{C}}
\providecommand{\e}{{\mathrm{e}}}
\providecommand{\vect}[1]{\overrightarrow{#1}}
\providecommand{\vectt}[1]{\overrightarrow{\mathrm{#1}}}
\providecommand{\vvt}[1]{\overrightarrow{\mathrm{#1}}}
\providecommand{\np}[2][]{{#2~#1}}
\providecommand{\vv}[1]{{\overrightarrow{#1}}}
\renewcommand{\d}{{\mathrm{d}}}
\providecommand{\degres}{{^\circ}}
\providecommand{\degre}{{^\circ}}
\providecommand{\blue}{\color{blue}}
\providecommand{\red}{\color{red}}
\providecommand{\enskip}{\quad}
\providecommand{\llbracket}{\lbrack\!\lbrack}
\providecommand{\rrbracket}{\rbrack\!\rbrack}
\providecommand{\euro}{{\text{€}}}
\providecommand{\no}{{\text{n°}}}
\providecommand{\interval}[2]{{[#1;#2]}}
\providecommand{\intervFO}[2]{{[#1;#2[}}
\providecommand{\dotfill}{\ldots}

\providecommand{\ds}{\displaystyle}
\providecommand{\textding}[1]{#1}
\providecommand{\Oij}{\(\mathrm{O};\vec{\imath},\vec{\jmath}\)}
\providecommand{\Oijk}{\(\mathrm{O};\vec{\imath},\vec{\jmath},\vec{k}\)}
\providecommand{\Ouv}{\(\mathrm{O};\vec{u},\vec{v}\)}
\providecommand{\intervFF}[2]{\ensuremath{[#1;#2]}}
\renewcommand{\euro}{\ifmmode\text{\texteuro}\else\texteuro\fi}
\DeclareUnicodeCharacter{00B0}{\ifmmode^{\circ}\else\textdegree\fi}
\DeclareUnicodeCharacter{20AC}{\ifmmode\text{\texteuro}\else\texteuro\fi}
% Énoncés : Éducation nationale. Transcription et corrigés : APMEP, Denis Vergès et contributeurs. Découpage et classement : CoopMaths, Sébastien Lozano, Éric Elter et contributeurs. Adaptation de la présentation : Maths Explorer.
% CC BY-SA : https://coopmaths.fr/www/about/
% https://coopmaths.fr/alea/static/bac/2025/tex/bac_2025_11_sujet2_caledonie_4.tex
\begin{document}
\noindent\textit{APMEP / CoopMaths. CC BY-SA. Présentation adaptée par Maths Explorer.}\par\medskip
On considère la fonction $f$ définie sur l'intervalle $]0~;~+\infty[$ par :

\[f(x) = \dfrac{\ln (x)}{x^{2}} + 1\]

On note $\mathcal{C}_{f}$ la courbe représentative de la fonction $f$ dans un repère orthonormé. On admet que la fonction $f$ est dérivable sur l'intervalle $]0~;~ +\infty[$ et on note $f'$ sa fonction dérivée.

\medskip

\begin{enumerate}
\item Déterminer les limites de la fonction $f$ en $0$ et en $+\infty$.

En déduire les éventuelles asymptotes à la courbe $\mathcal{C}_{f}$.
\item Montrer que, pour tout réel $x$ de l'intervalle $]0~;~ +\infty[$, on a :

\[f'(x) = \dfrac{1 - 2\ln(x)}{x^{3}}\]

\item En déduire le tableau de variation de la fonction $f$ sur l'intervalle $]0~;~ +\infty[$.

\item 
	\begin{enumerate}
		\item Montrer que l'équation $f(x) = 0$ possède une unique solution, notée $\alpha$, sur l'intervalle $]0~;~ +\infty[$.
		\item Donner un encadrement du réel $\alpha$ d'amplitude $0,01$.
		\item En déduire le signe de la fonction $f$ sur l'intervalle $]0~;~ +\infty[$.
	\end{enumerate}
\item On considère la fonction $g$ définie sur l'intervalle $]0~;~ +\infty[$ par :

\[g(x) = \ln(x)\]

On note $\mathcal{C}_{g}$ la courbe représentative de la fonction $g$ dans un repère orthonormé d'origine O. On considère un réel $x$ strictement positif et le point M de la courbe $C_{g}$ d'abscisse $x$. On note OM la distance entre les points O et M.

	\begin{enumerate}
		\item Exprimer la quantité OM$^{2}$ en fonction du réel $x$.
		\item Montrer que, lorsque le réel $x$ parcourt l'intervalle $]0~;~ +\infty[$, la quantité OM$^{2}$ admet un minimum en $\alpha$.
		\item La valeur minimale de la distance OM, lorsque le réel $x$ parcourt l'intervalle $]0~;~+\infty[$, est appelée distance du point O à la courbe $\mathcal{C}_{g}$. On note $d$ cette distance.

Exprimer $d$ à l'aide de $\alpha$.
	\end{enumerate}
\end{enumerate}

\end{document}
